% Vietnamese translation for OI Public Library
% Source-File: IOI中国国家候选队论文/2006/黄劲松/黄劲松.ppt
\documentclass[11pt]{article}
\usepackage{oipl-vn}

\newcommand{\Oh}{\mathcal{O}}

\title{Quy hoạch động tham lam\\{\large Ghi chú theo slide}}
\author{Huang Jinsong}
\date{2006}

\begin{document}
\maketitle

\origtitle{Greedy dynamic programming}
\sourcepath{IOI China candidate team papers / 2006 / Huang Jinsong / Huang Jinsong.ppt}

\section{Phạm vi}

Bộ trình chiếu đi kèm bài viết về việc đưa tư tưởng tham lam vào thiết kế và
tối ưu quy hoạch động. Ghi chú này dựa trên bản bài viết đã dịch.

\section{Tham lam không thay thế quy hoạch động}

Trong các bài phù hợp với quy hoạch động, tham lam thường không tự giải toàn
bộ bài toán. Vai trò của nó là tìm tính chất cực trị của nghiệm, từ đó giảm
lượng thông tin cần lưu trong trạng thái hoặc loại bỏ các chuyển chắc chắn
không thể tối ưu.

\section{Xác lập trạng thái}

Ví dụ chú ếch trên đa giác lồi dùng quan sát tham lam kiểu 2-opt: hành trình
tối ưu không có hai cạnh cắt nhau. Từ đó trạng thái chỉ cần là một đoạn liên
tiếp của các đỉnh và vị trí ở một trong hai đầu đoạn, thay vì một tập đỉnh bất
kỳ như bài TSP tổng quát.

\section{Tối ưu chuyển trạng thái}

Ở nhóm bài thứ hai, công thức quy hoạch động ban đầu đúng nhưng có quá nhiều
quyết định. Một chứng minh tham lam có thể cho thấy quyết định tốt nhất có
tính đơn điệu, hoặc một số quyết định bị quyết định khác lấn át. Khi đó ta
dùng hàng đợi, chia để trị, hoặc cấu trúc dữ liệu phù hợp để giảm thời gian.

\section{Ghi nhớ}

Hãy tìm phần của nghiệm tối ưu có thể khẳng định chắc chắn trước khi viết
trạng thái. Một mệnh đề tham lam được chứng minh tốt thường là chìa khóa làm
cho quy hoạch động vừa nhỏ vừa dễ cài đặt.

\end{document}
